Дано целое число. Если оно является положительным и четным, то прибавить к нему 1; в противном случае не изменять его. Вывести полученное число.
Дано целое число. Если оно является положительным и нечетным, то прибавить к нему 1; в противном случае вычесть из него 2. Вывести полученное число.
Дано целое число. Если оно является положительным, то прибавить к нему 1; если отрицательным, то вычесть из него 2; если нулевым, то заменить его на 10. Вывести полученное число.
Даны три целых числа. Найти количество отрицательных и количество четных чисел в исходном наборе.
Даны три целых числа. Найти количество положительных и количество отрицательных чисел в исходном наборе.
Даны два числа. Вывести большее из них.
Даны два числа. Вывести порядковый номер меньшего из них.
Даны два числа. Вывести вначале большее, а затем меньшее из них.
Даны две переменные вещественного типа: $A$, $B$. Перераспределить значения данных переменных так, чтобы в $A$ оказалось меньшее из значений, а в $B$ --- большее. Вывести новые значения переменных $A$ и $B$.
Даны две переменные целого типа: $A$ и $B$. Если их значения не равны, то присвоить каждой переменной сумму этих значений, а если равны, то присвоить переменным нулевые значения. Вывести новые значения переменных $A$ и $B$.
Даны две переменные целого типа: $A$ и $B$. Если их значения не равны, то присвоить каждой переменной большее из этих значений, а если равны, то присвоить переменным нулевые значения. Вывести новые значения переменных $A$ и $B$.
Даны три числа. Найти наименьшее из них.
Даны три числа. Найти среднее из них (то есть число, расположенное между наименьшим и наибольшим).
Даны три числа. Вывести вначале наименьшее, а затем наибольшее из данных чисел.
Даны три числа. Найти сумму двух наибольших из них.
Даны три переменные вещественного типа: $A$, $B$, $C$. Если их значения упорядочены по возрастанию, то удвоить их; в противном случае заменить значение каждой переменной на противоположное. Вывести новые значения переменных $A$, $B$, $C$.
Даны три переменные вещественного типа: $A$, $B$, $C$. Если их значения упорядочены по возрастанию или убыванию, то удвоить их; в противном случае заменить значение каждой переменной на противоположное. Вывести новые значения переменных $A$, $B$, $C$.
Даны три целых числа, одно из которых отлично от двух других, равных между собой. Определить порядковый номер числа, отличного от остальных.
Даны четыре целых числа, одно из которых отлично от трех других, равных между собой. Определить порядковый номер числа, отличного от остальных.
На числовой оси расположены три точки: $A$, $B$, $C$. Определить, какая из двух последних точек ($B$ или $C$) расположена ближе к $A$, и вывести эту точку и ее расстояние от точки $A$.
Даны целочисленные координаты точки на плоскости. Если точка совпадает с началом координат, то вывести $0$. Если точка не совпадает с началом координат, но лежит на оси $OX$ или $OY$, то вывести соответственно $1$ или $2$. Если точка не лежит на координатных осях, то вывести $3$.
Даны координаты точки, не лежащей на координатных осях $OX$ и $OY$. Определить номер координатной четверти, в которой находится данная точка.
Даны целочисленные координаты трех вершин прямоугольника, стороны которого параллельны координатным осям. Найти координаты его четвертой вершины.
Дано целое число. Вывести его строку-описание вида «отрицательное четное число», «нулевое число», «положительное нечетное число» и т. д.
Дано целое число, лежащее в диапазоне $1$--$999$. Вывести его строку-описание вида «четное двузначное число», «нечетное трехзначное число» и т. д.
Даны три целых числа. Найти количество положительных и количество нечетных чисел в исходном наборе.
Даны три целых числа. Найти количество четных и количество нечетных чисел в исходном наборе.
Даны три числа. Найти произведение двух наименьших из них.
Дано целое число, лежащее в диапазоне $1$--$999$. Вывести его строку-описание вида «положительное двузначное число», «отрицательное трехзначное число» и т. д.
Даны три числа. Найти произведение двух наибольших из них.
